\newpage
\phantomsection
\label{prf:q-is-integral-domain}

\begin{remark*}[Return]
\hyperref[thm:q-is-integral-domain]{Return to Theorem}
\end{remark*}

\begin{theorem*}[$\mathbb{Q}$ Is an Integral Domain]
With the operations of
\hyperref[def:rational-operations]{Definition~\ref*{def:rational-operations}},
the system
\[
(\mathbb{Q},+,\cdot,0,1)
\]
is an integral domain.
\end{theorem*}

\begin{proof}
\textbf{Professional Standard Proof.}~\\
We verify the five conditions of the integral domain definition.

\textbf{(1) $(\mathbb{Q},+,0)$ is a commutative group.}

\emph{Closure.} If $[(a,b)],[(c,d)]\in\mathbb{Q}$ then
$[(ad+bc,bd)]\in\mathbb{Q}$ since $bd\neq 0$.

\emph{Associativity.} Let $p = [(a,b)]$, $q = [(c,d)]$, $r = [(e,f)]$.
Then
\begin{align*}
(p+q)+r &= [(ad+bc,bd)] + [(e,f)] = [((ad+bc)f + bd\cdot e,\; bdf)], \\
p+(q+r) &= [(a,b)] + [(cf+de,df)] = [(a\cdot df + b(cf+de),\; bdf)].
\end{align*}
The numerators are $(adf + bcf + bde)$ in both cases, so the sums are equal.

\emph{Identity.} $[(a,b)] + [(0,1)] = [(a\cdot 1 + b\cdot 0, b\cdot 1)] = [(a,b)]$.

\emph{Inverses.} $[(a,b)] + [(-a,b)] = [(ab - ab, b^2)] = [(0,b^2)] = 0$,
since $(0,b^2)\sim(0,1)$.

\emph{Commutativity.} $[(a,b)]+[(c,d)] = [(ad+bc,bd)] = [(cb+da,db)] = [(c,d)]+[(a,b)]$.

\textbf{(2) Multiplication is associative and commutative.}

\emph{Associativity.}
$([(a,b)]\cdot[(c,d)])\cdot[(e,f)] = [(ac,bd)]\cdot[(e,f)] = [(ace,bdf)]$.
$[(a,b)]\cdot([(c,d)]\cdot[(e,f)]) = [(a,b)]\cdot[(ce,df)] = [(ace,bdf)]$.

\emph{Commutativity.}
$[(a,b)]\cdot[(c,d)] = [(ac,bd)] = [(ca,db)] = [(c,d)]\cdot[(a,b)]$.

\textbf{(3) $1\neq 0$ and $1$ is the multiplicative identity.}

$[(1,1)]\neq[(0,1)]$ since $1\cdot 1 = 1\neq 0 = 0\cdot 1$.
$[(1,1)]\cdot[(a,b)] = [(1\cdot a, 1\cdot b)] = [(a,b)]$.

\textbf{(4) Multiplication distributes over addition.}

\begin{align*}
[(a,b)]\cdot([(c,d)]+[(e,f)])
&= [(a,b)]\cdot[(cf+de,df)] \\
&= [(a(cf+de),\; b\cdot df)] \\
&= [(acf+ade,\; bdf)].
\end{align*}
\begin{align*}
[(a,b)]\cdot[(c,d)] + [(a,b)]\cdot[(e,f)]
&= [(ac,bd)] + [(ae,bf)] \\
&= [(ac\cdot bf + bd\cdot ae,\; bd\cdot bf)] \\
&= [(abcf + abde,\; b^2df)].
\end{align*}
These two classes are equal iff $(acf+ade)(b^2df) = (abcf+abde)(bdf)$,
which simplifies to $ab^2df(cf+de) = ab(bdf)(cf+de)$, i.e.\ $b^2df = b^2df$.
True.

\textbf{(5) $\mathbb{Q}$ has no zero divisors.}

Suppose $[(a,b)]\cdot[(c,d)] = 0$. Then $[(ac,bd)] = [(0,1)]$, so
$(ac,bd)\sim(0,1)$, meaning $ac\cdot 1 = bd\cdot 0 = 0$. Thus $ac = 0$.
Since $\mathbb{Z}$ is an integral domain, $a=0$ or $c=0$. By the zero class
criterion, $a=0$ iff $[(a,b)] = 0$ and $c=0$ iff $[(c,d)] = 0$.
Therefore $[(a,b)]=0$ or $[(c,d)]=0$.

Since all five conditions hold, $(\mathbb{Q},+,\cdot,0,1)$ is an integral domain.
\end{proof}

\begin{proof}
\textbf{Detailed Learning Proof.}~\\

We check each condition of the integral domain definition in turn.

\textbf{Condition 1: $(\mathbb{Q},+,0)$ is a commutative group.}

We must verify five sub-properties.

\emph{(1a) Closure under addition.}
Let $[(a,b)],[(c,d)]\in\mathbb{Q}$. The sum is $[(ad+bc,\,bd)]$. Since
$b\neq 0$ and $d\neq 0$, we have $bd\neq 0$, so $(ad+bc,bd)\in W$ and
$[(ad+bc,bd)]\in\mathbb{Q}$.

\emph{(1b) Associativity of addition.}
Let $[(a,b)],[(c,d)],[(e,f)]\in\mathbb{Q}$. We compute both bracketings
directly:
\[
([(a,b)]+[(c,d)])+[(e,f)]
= [(ad+bc,bd)]+[(e,f)]
= [((ad+bc)f+bd\cdot e,\; bdf)]
= [(adf+bcf+bde,\; bdf)].
\]
\[
[(a,b)]+([(c,d)]+[(e,f)])
= [(a,b)]+[(cf+de,df)]
= [(a\cdot df + b(cf+de),\; bdf)]
= [(adf+bcf+bde,\; bdf)].
\]
The numerators and denominators are identical, so
$([(a,b)]+[(c,d)])+[(e,f)] = [(a,b)]+([(c,d)]+[(e,f)])$.

\emph{(1c) Additive identity.}
$0 = [(0,1)]$. Compute:
\[
[(a,b)] + [(0,1)]
= [(a\cdot 1 + b\cdot 0,\; b\cdot 1)]
= [(a,b)].
\]
Similarly $[(0,1)] + [(a,b)] = [(a,b)]$ by commutativity of integer addition.

\emph{(1d) Additive inverses.}
The candidate inverse of $[(a,b)]$ is $[(-a,b)]$ (see the negation theorem).
We have already verified this: $[(a,b)]+[(-a,b)] = [(0,b^2)] = 0$.

\emph{(1e) Commutativity of addition.}
\[
[(a,b)]+[(c,d)] = [(ad+bc,bd)] = [(cb+da,db)] = [(c,d)]+[(a,b)],
\]
where we used commutativity of integer addition ($ad+bc = cb+da$) and
commutativity of integer multiplication ($bd = db$).

\textbf{Condition 2: Multiplication is associative and commutative.}

\emph{(2a) Associativity.}
\[
([(a,b)]\cdot[(c,d)])\cdot[(e,f)]
= [(ac,bd)]\cdot[(e,f)]
= [(ace,bdf)].
\]
\[
[(a,b)]\cdot([(c,d)]\cdot[(e,f)])
= [(a,b)]\cdot[(ce,df)]
= [(ace,bdf)].
\]
Both equal $[(ace,bdf)]$, so multiplication is associative.

\emph{(2b) Commutativity.}
$[(a,b)]\cdot[(c,d)] = [(ac,bd)] = [(ca,db)] = [(c,d)]\cdot[(a,b)]$,
using commutativity in $\mathbb{Z}$.

\textbf{Condition 3: $1\neq 0$ and $1$ is the multiplicative identity.}

$1 = [(1,1)]$ and $0 = [(0,1)]$. These are not equal: $(1,1)\sim(0,1)$
would require $1\cdot 1 = 0\cdot 1$, i.e.\ $1 = 0$, which is false in
$\mathbb{Z}$. So $1\neq 0$.

For the identity property:
\[
[(1,1)]\cdot[(a,b)]
= [(1\cdot a,\; 1\cdot b)]
= [(a,b)].
\]
Commutativity gives $[(a,b)]\cdot[(1,1)] = [(a,b)]$ as well.

\textbf{Condition 4: Distributivity.}

We show $[(a,b)]\cdot([(c,d)]+[(e,f)]) = [(a,b)]\cdot[(c,d)] + [(a,b)]\cdot[(e,f)]$.

Left-hand side:
\[
[(a,b)]\cdot[(cf+de,\;df)]
= [(a(cf+de),\;bdf)]
= [(acf+ade,\;bdf)].
\]
Right-hand side:
\[
[(ac,bd)] + [(ae,bf)]
= [(ac\cdot bf + bd\cdot ae,\; bd\cdot bf)]
= [(abcf+abde,\; b^2df)].
\]
We must check $[(acf+ade,bdf)] = [(abcf+abde,b^2df)]$.
By the definition of $\sim$:
\[
(acf+ade)(b^2df) \stackrel{?}{=} (bdf)(abcf+abde).
\]
Left: $b^2df(acf+ade) = ab^2df(cf+de)$.\\
Right: $bdf\cdot ab(cf+de) = ab^2df(cf+de)$.\\
These are equal. Therefore distributivity holds.

\textbf{Condition 5: No zero divisors.}

Let $[(a,b)],[(c,d)]\in\mathbb{Q}$ and suppose
\[
[(a,b)]\cdot[(c,d)] = 0.
\]
By the multiplication formula, $[(ac,bd)] = [(0,1)]$.
By the zero class criterion applied to $[(ac,bd)]$, this means $ac = 0$.
Since $\mathbb{Z}$ is an integral domain (integers have no zero divisors),
$ac = 0$ implies $a = 0$ or $c = 0$.

If $a = 0$, then by the zero class criterion $[(a,b)] = 0$.\\
If $c = 0$, then by the zero class criterion $[(c,d)] = 0$.

Therefore at least one of the factors is zero in $\mathbb{Q}$, and $\mathbb{Q}$
has no zero divisors.

\textbf{Conclusion.}

All five conditions of the integral domain definition are satisfied.
Therefore $(\mathbb{Q},+,\cdot,0,1)$ is an integral domain.
\end{proof}

\begin{remark*}[Proof structure]
The proof is a systematic verification of each axiom. The substantive
content lies in three places: the associativity of addition (which requires
expanding both sides and comparing numerators), the distributivity check
(which reduces to a cross-multiplication equality), and the no-zero-divisors
property (which reduces $\mathbb{Q}$-level zeroes to the corresponding fact
in $\mathbb{Z}$, which is an integral domain by construction).

All other axioms are direct computations from the operation definitions and
the arithmetic of $\mathbb{Z}$.
\end{remark*}

\begin{remark*}[Dependencies]~\\
\begin{itemize}
  \item \hyperref[def:rational-numbers]{Definition~\ref*{def:rational-numbers}}
  \item \hyperref[def:rational-distinguished]{Definition~\ref*{def:rational-distinguished}}
  \item \hyperref[def:rational-operations]{Definition~\ref*{def:rational-operations}}
  \item \hyperref[def:integral-domain]{Definition~\ref*{def:integral-domain}}
  \item \hyperref[thm:ec-zero-numerator]{Theorem~\ref*{thm:ec-zero-numerator}}
        \textnormal{(zero class criterion, used in Condition 5)}
  \item $\mathbb{Z}$ is an integral domain (no zero divisors in $\mathbb{Z}$)
\end{itemize}
\end{remark*}
